\section{Referencial teórico}

A comensuração permite examinar como diferenças são traduzidas em medidas comparáveis, sem tratar essa tradução como operação socialmente neutra \cite{espeland1998comensuracao}. Essa lente fundamenta a análise de classificações, denominadores e formas de agregação dos indicadores de violência letal. Em outra arena, \citeonline{kristensen2023indicadores} investigam relações entre indicadores de problemas e atenção parlamentar em sete países. O estudo sustenta a pertinência de observar a mobilização política de indicadores, mas não constitui evidência direta sobre o Brasil nem demonstra o percurso documental que esta pesquisa pretende reconstruir.

\citeonline{pilar2020atencao} comparam causas de morte, atenção midiática, atenção legislativa e financiamento nos Estados Unidos, indicando que as posições relativas nessas arenas não são necessariamente equivalentes. A contribuição para este projeto é analítica: comparar arenas exige preservar seus objetos, unidades e métricas. O esqueleto mantém a delimitação documental já explicitada pelo autor no pré-projeto. A transcrição a seguir é de um documento autoral fornecido, não de resultado obtido neste exercício, e demonstra a citação direta longa.

\begin{citacao}
O estudo buscará reconstruir um mecanismo documental de circulação, seleção, reconfiguração e mobilização, a partir de rastros empiricamente verificáveis entre as arenas. Uma citação explícita de uma fonte institucional constitui evidência mais forte de ligação documental; a reprodução do mesmo indicador, valor ou comparação constitui evidência de correspondência; a mera sucessão temporal, por sua vez, será tratada apenas como contexto, e não como demonstração de transmissão ou causalidade. \cite[p.~10]{santos2026preprojeto}
\end{citacao}
